Ultrasound fibrosis grading. Prior work has used B-mode ultrasound, elastography, micro-flow measurements, and generative augmentation for fibrosis assessment~\cite{microflow2023,swe2018staging,usfibrosisgan2022}. These approaches generally target a small number of stages and often focus on etiologies other than schistosomiasis. SFibAI introduced a 36-level posterior-expectation formulation for fine-grained grading in this setting~\cite{xu2026sfibai}. We retain that formulation and use it as the primary task on which to study concept pathways.

Concept-based and interpretable prediction. Concept Bottleneck Models predict human-specified concepts before a final target and support interventions on those concepts~\cite{koh20a}. Related medical-imaging systems use segmentation, anatomy-supervised attention, or auxiliary heads~\cite{anatomyxnet2022,zhou2020mtlabus}, but commonly leave the concept-to-task relationship entangled with shared gradients. Our model is a non-bottleneck concept pathway: the grading backbone remains available, while predicted anatomical concepts enter through explicit residual interfaces. This distinction permits concept ablations and future counterfactual interventions without claiming that the current model is a strict bottleneck.

Anatomy-informed auxiliary learning. Anatomy-guided networks and multi-task models have used organ masks, poses, parts, and weak localization to shape visual representations. SynAP-Fib studies a different setting in which the anatomical signals are heterogeneous: position is categorical and lesion evidence is spatial and weakly labelled. The same interface handles both signals, and the $2\times2$ comparison asks whether their joint use behaves differently from either signal alone.

Positioning. The novelty claim is deliberately narrow. We do not introduce a new ordinal loss, a new backbone, or a universal explanation algorithm. We provide a controlled medical-vision testbed for image-derived anatomical concepts, a gradient-isolated concept interface, and a complete comparison of single versus joint concept pathways. Direct intervention tests are the next step required to establish faithfulness.


